% Личный материал (не часть пакета most-queue, docs/статьи/ в .gitignore).
% Стиль оформления -- по образцу реальной опубликованной статьи автора в
% журнале "Труды Военно-космической академии имени А.Ф. Можайского", тот же
% шаблон, что в batch-service-sla-single/multi.
%
% Сборка: latexmk -xelatex main.tex   (требуется XeLaTeX + biber)
%
% encoding: UTF-8
\documentclass[12pt,a4paper]{article}

\usepackage{geometry}
\geometry{top=2.5cm,bottom=2.5cm,left=3cm,right=2cm}

\usepackage{fontspec}
\setmainfont{PT Serif}
\newfontfamily\cyrillicfonttt{DejaVu Sans Mono}

\usepackage{polyglossia}
\setmainlanguage{russian}
\setotherlanguage[variant=american]{english}

\usepackage{amsmath,amssymb}
\usepackage{graphicx}
\usepackage{siunitx}
\sisetup{output-decimal-marker = {,}, group-separator = {\,}, retain-unity-mantissa = false}

\usepackage{array}
\usepackage{multirow}
\usepackage{caption}
\usepackage{capt-of}
\usepackage{seqsplit}
\usepackage{xcolor}
\usepackage{tikz}
\usetikzlibrary{arrows.meta, positioning, shapes.geometric}
\captionsetup[figure]{labelsep=period, justification=centering, font=small, position=bottom}
\captionsetup[table]{labelsep=period, font={small,bf}, justification=raggedright, position=top}

\usepackage{indentfirst}
\setlength{\parindent}{1.25cm}
\linespread{1.12}

\usepackage[backend=biber, style=gost-numeric, sorting=none]{biblatex}
\addbibresource{../../bibliography/library.bib}

% Заголовки разделов: без нумерации, жирным, без капслока -- по образцу
% реальной статьи автора в "Трудах ВКА" (не ГОСТовский численный рубрикатор).
\newcommand{\articlesection}[1]{\par\vspace{14pt}{\bfseries #1}\par\vspace{6pt}\nopagebreak}

\newcommand{\keywordslabel}{\textbf{Ключевые слова: }}

\begin{document}
\sloppy

\noindent УДК 519.872

\bigskip
\noindent Р.С.~Хабаров, к.т.н.

\bigskip
\begin{center}
\begin{minipage}{0.9\linewidth}
\centering
{\large\bfseries ТОЧНОЕ РАСПРЕДЕЛЕНИЕ ВРЕМЕНИ ОЖИДАНИЯ В СИСТЕМЕ МАССОВОГО ОБСЛУЖИВАНИЯ С OCCUPANCY-ЗАВИСИМОЙ ИНТЕНСИВНОСТЬЮ И ПОТОЛКОМ ЗАНЯТОСТИ}
\end{minipage}
\end{center}

\bigskip
\noindent\small
Рассмотрена одноканальная (в смысле общего разделяемого ресурса) система массового обслуживания, в которой до $k$ заявок обслуживаются одновременно, интенсивность обслуживания каждой заявки зависит от текущего числа одновременно обслуживаемых (occupancy-зависимая интенсивность обслуживания), а превышающие потолок $k$ заявки ожидают по дисциплине FCFS. Такая система -- точная марковская модель ядра <<continuous batching>> (непрерывного, поитерационного формирования батча) современных движков инференса больших языковых моделей (LLM), где $k$ задаётся объёмом доступной памяти ускорителя (KV-cache), а occupancy-зависимость интенсивности обслуживания отражает разделяемую вычислительную мощность/пропускную способность памяти между одновременно обслуживаемыми запросами. Показано, что стационарное распределение числа заявок в системе имеет точную (не итеративную) мультипликативную форму в силу детального баланса, присущего любому процессу гибели-размножения; показано также, что занятость канала жёстко зафиксирована на уровне $k$ в любом состоянии с непустой очередью, откуда следует точное представление времени ожидания PASTA-наблюдающей заявки в виде явной смеси эрланговских распределений -- точные начальные моменты и точный хвост $P(W>D)$, а не только среднее значение, как в единственной найденной близкой по постановке работе. Результат точно регрессирует к классической системе $M/M/k/N$ при occupancy-независимой интенсивности и к $M/M/1/N$ при $k=1$; подтверждён независимым имитационным моделированием. Численно показано, что простое приближение, игнорирующее occupancy-зависимость интенсивности (использующее интенсивность, измеренную при полной занятости, как константу), систематически завышает вероятность нарушения дедлайна -- на один-два порядка при умеренной занятости и на несколько порядков при большом потолке занятости $k$.

\medskip
\noindent\keywordslabel теория массового обслуживания, occupancy-зависимое обслуживание, процесс гибели-размножения, потолок занятости, continuous batching, инференс больших языковых моделей, KV-cache, точное распределение времени ожидания.
\normalsize

\bigskip
\begin{center}
\begin{minipage}{0.9\linewidth}
\centering
{\bfseries EXACT WAITING-TIME DISTRIBUTION FOR AN OCCUPANCY-DEPENDENT QUEUE WITH A CONCURRENCY CAP: THE MARKOVIAN CORE OF LLM-SERVING CONTINUOUS BATCHING}
\end{minipage}
\end{center}

\begin{center}
\small R.S.~Khabarov
\end{center}

\begin{otherlanguage}{english}
\small
We consider a single shared-resource queueing system in which up to $k$ customers are served concurrently, the per-customer completion rate depends on the current number of concurrently-served customers (an occupancy-dependent rate), and customers beyond the cap $k$ wait FCFS. This system is an exact Markovian model of the core mechanism of ``continuous batching'' in modern large language model (LLM) inference engines, where $k$ is set by available accelerator memory (the KV-cache) and the occupancy-dependent rate reflects shared compute/memory-bandwidth capacity among concurrently-resident requests. We show that the stationary distribution of the number in system has an exact (non-iterative) product form, a consequence of detailed balance for any birth-death process; we further show that occupancy is pinned at exactly $k$ in every state with a nonempty queue, from which an exact representation of the waiting time of a PASTA-observing customer follows as an explicit mixture of Erlang distributions -- exact raw moments and an exact tail $P(W>D)$, not just the mean, unlike the one directly comparable prior treatment of this system class. The result regresses exactly to the classical $M/M/k/N$ queue at an occupancy-independent rate and to $M/M/1/N$ at $k=1$; it is confirmed against independent discrete-event simulation. Numerically, a simple approximation that ignores the occupancy dependence (using the rate measured at full occupancy as a constant) is shown to systematically overstate deadline-violation probability -- by one to two orders of magnitude at moderate occupancy and by several orders of magnitude at a large concurrency cap $k$.

\medskip
\noindent\textit{Key words:} queueing theory, occupancy-dependent service, birth-death process, concurrency cap, continuous batching, large language model inference, KV-cache, exact waiting-time distribution.
\end{otherlanguage}
\normalsize

\renewcommand{\figurename}{Рисунок}
\renewcommand{\tablename}{Таблица}

\articlesection{Введение}

Современные движки инференса больших языковых моделей (vLLM, SGLang и близкие по архитектуре) используют <<непрерывный>> (continuous, или поитерационный) батчинг: в отличие от классического пакетного обслуживания \cite{bailey1954}, где канал копит от $a$ до $b$ заявок и запускает обслуживание всего пакета разом \cite{inoue2021}, здесь состав одновременно обслуживаемых запросов меняется на каждой итерации декодирования -- запрос присоединяется к уже исполняемой группе, как только освобождается место, и покидает её независимо от остальных. Число мест, которые могут быть заняты одновременно ($k$), задаётся не произвольным порогом, а объёмом доступной памяти ускорителя для KV-cache, а интенсивность обслуживания каждого запроса, как правило, зависит от текущей занятости: все резидентные запросы делят одну итерацию вычислений, и большая группа означает большую амортизацию, но и большую конкуренцию за пропускную способность памяти.

Такая occupancy-зависимая интенсивность обслуживания с потолком занятости структурно отличается от уже рассмотренного в \cite{inoue2021} и в работах автора \cite{khabarovsingle2026,khabarovmulti2026,khabarovoccupancy2026} атомарного пакетного обслуживания: здесь нет понятия <<пакет>>, завершающегося целиком, а есть обычный процесс гибели-размножения по числу заявок в системе, но с зависящей от состояния интенсивностью смерти.

Системы массового обслуживания с occupancy-зависимой (state-dependent) интенсивностью обслуживания известны с 1970-х годов \cite{morrison1989} и ранее -- сама по себе такая цепь не нова. Точное распределение времени ожидания для неё, однако, найдено (насколько можно судить по проведённому обзору) лишь для одно- и двухсерверного случаев с двухуровневой (пороговой) интенсивностью \cite{morrison1989}. Специализированная монография по этой теме \cite{bekker2005} (защищена под руководством одного из ведущих авторов данного направления) целиком ограничена однoсерверным случаем ($G/G/1$-постановки: зависимость от объёма работы, дамбы и накопители, оптимальный допуск, адаптируемая скорость обслуживания) -- ни одна из её семи глав не рассматривает многоканальную систему с потолком одновременной занятости. Обобщение на произвольное число одновременно обслуживаемых заявок $k$ и произвольный профиль интенсивности по occupancy в явном виде не найдено.

Современная литература по LLM-serving активно рассматривает ту же систему применительно к continuous batching, но ограничивается либо приближённым, либо неполным результатом. Yang, Xu и Jiao \cite{yang2024} формулируют M/G/1 и пакетную модель с учётом длины вывода, но получают только среднее время ожидания и не рассматривают occupancy-зависимую интенсивность с потолком занятости как таковую. Ramani и Tantawi \cite{ramani2026} строят ту же <<state-dependent Markov chain for batch occupancy>>, что и в настоящей работе, но решают её приближённо, методом среднего значения (mean-value analysis), и дают только среднее время до первого токена и средний межтокенный интервал (ошибка валидации против реальных замеров -- 5--16\,\%). Nie, Si и Zhou \cite{nie2026} дают условия устойчивости/неустойчивости такой системы, но не стационарное распределение времени ожидания. Ao с соавторами \cite{ao2026} рассматривают смежный, но принципиально иной вопрос -- детерминированную (не стохастическую) динамику перегрузки при вытеснении KV-cache, предельные циклы и потерю пропускной способности, другим математическим аппаратом (динамические системы, не цепи Маркова).

Таким образом, пересечение <<точно>>$\times$<<полное распределение (не только среднее)>>$\times$<<произвольное $k$>>$\times$<<произвольный профиль occupancy-зависимой интенсивности>> не закрыто ни одной из найденных работ. Настоящая статья даёт такое точное решение, включая точный хвост $P(W>D)$, представляющий прямой практический интерес для SLA-ориентированного управления инференс-серверами.

\articlesection{Постановка задачи}

Рассматривается система с пуассоновским входным потоком интенсивности $\lambda$, в которой одновременно может обслуживаться не более $k$ заявок; интенсивность обслуживания каждой из них, пока занято ровно $n$ мест ($1\le n\le k$), равна $\mu(n)$ -- общая, не обязательно постоянная зависимость от occupancy $n$ (в частности, occupancy-независимая интенсивность -- частный случай). Заявки, заставшие все $k$ мест занятыми, ожидают по FCFS; как только освобождается место, его немедленно занимает первая заявка из очереди (если очередь не пуста).

Цель -- точное распределение времени ожидания $W$ PASTA-наблюдающей (случайно выбранной) заявки до момента её допуска в число одновременно обслуживаемых, при произвольном $k$ и произвольной зависимости $\mu(n)$, $1\le n\le k$.

\articlesection{Метод}

\noindent\textbf{Стационарное распределение числа в системе -- известный результат (метод Эрланга).} Пусть $j$ -- полное число заявок в системе (обслуживаемых и ожидающих), $0\le j\le N$ (усечение по $N$ для вычислений). Занятость в состоянии $j$ равна $\min(j,k)$, поэтому суммарная интенсивность выхода из состояния $j$ (<<интенсивность смерти>>) равна
\begin{equation}
\delta(j) = \min(j,k)\cdot\mu(\min(j,k)).
\label{eq:death-rate}
\end{equation}
Цепь по $j$ -- процесс гибели-размножения с интенсивностью рождения $\lambda$ и интенсивностью смерти $\delta(j)$. Для произвольного процесса гибели-размножения со сколь угодно зависящими от состояния интенсивностями стационарное распределение в замкнутой, не итеративной форме -- классический результат, известный как метод Эрланга \cite{ryzhikov2018}:
\begin{equation}
\pi_j = \pi_0\prod_{i=1}^{j}\frac{\lambda}{\delta(i)}, \qquad \pi_0 = \left(\sum_{j=0}^{N}\prod_{i=1}^{j}\frac{\lambda}{\delta(i)}\right)^{-1}.
\label{eq:pi-product-form}
\end{equation}
Там же выписан частный случай $n$-линейной системы $M/M/n$, для которого при $j>n$ интенсивность смерти не зависит от $j$ и равна $n\mu$ -- тот же классический факт, на котором строится формула Эрланга-C, нужен нам в следующем пункте уже для произвольного (не обязательно постоянного ниже потолка) профиля $\mu(n)$.

\noindent\textbf{Интенсивность смерти выше потолка -- тот же классический факт, перенесённый на произвольный профиль.} Для всех $j\ge k$ занятость равна в точности $k$, поэтому из формулы~\eqref{eq:death-rate} непосредственно $\delta(j)=k\cdot\mu(k)$ -- не зависит от $j$, какой бы ни была зависимость $\mu(n)$ при $n<k$ (это следует из построения: пока очередь не пуста, любое освобождение места немедленно восполняется из очереди, поэтому занятость не может опуститься ниже $k$). Новое здесь -- не сам этот факт (он тот же, что и в формуле Эрланга-C выше), а то, что он верен при произвольном профиле $\mu(n)$ ниже потолка, а не только при постоянном: ниже мы используем это, чтобы получить точное распределение времени ожидания, а не только стационарные вероятности состояний.

\noindent\textbf{Время ожидания.} По свойству PASTA заявка, приходящая в систему, видит стационарное распределение $\pi_j$. Если $j<k$ -- заявка допускается немедленно, $W=0$. Если $j\ge k$ -- заявка становится $(j-k+1)$-й в очереди и должна дождаться ровно $(j-k+1)$ освобождений мест; эти освобождения происходят с постоянной интенсивностью $k\cdot\mu(k)$ независимо от того, что происходит ниже потолка -- то есть условное (на $j\ge k$) время ожидания есть точный $\mathrm{Erlang}(j-k+1,\,k\cdot\mu(k))$, ровно как для тех же $(j-k+1)$ заявок впереди в классической системе $M/M/k$. Безусловные точные начальные моменты и точный хвост:
\begin{equation}
\mathrm{E}[W^n] = \sum_{j=k}^{N}\pi_j\,\frac{(j-k+1)(j-k+2)\cdots(j-k+n)}{(k\,\mu(k))^n},
\label{eq:w-moments}
\end{equation}
\begin{equation}
P(W>t) = \sum_{j=k}^{N}\pi_j\sum_{i=0}^{j-k}\frac{e^{-k\,\mu(k)\,t}\,(k\,\mu(k)\,t)^i}{i!}.
\label{eq:w-tail}
\end{equation}

На рисунке~\ref{fig:birth-death} схематично показана цепь: состояния $j<k$ образуют occupancy-зависимую зону (интенсивность зависит от $j$), состояния $j\ge k$ -- зону с постоянной, <<прижатой>> к потолку интенсивностью.

\begin{center}
\resizebox{0.95\linewidth}{!}{%
\begin{tikzpicture}[
  every node/.style={font=\footnotesize},
  state/.style={ellipse, draw=black, thick, minimum width=11mm, minimum height=8mm, inner sep=0pt},
  var/.style={state, fill=green!18},
  pinned/.style={state, fill=blue!15},
  arr/.style={-{Stealth[length=2.2mm]}, thick, blue!55!black},
  dep/.style={-{Stealth[length=2.2mm]}, thick, purple!65!black},
]
\node[var] (s0) at (0,0) {$0$};
\node[var] (s1) at (2.2,0) {$1$};
\node[var] (s2) at (4.4,0) {$2$};
\node at (6.2,0) {$\cdots$};
\node[var] (sk1) at (8.0,0) {$k{-}1$};
\node[pinned] (sk) at (10.6,0) {$k$};
\node[pinned] (sk2) at (13.2,0) {$k{+}1$};
\node at (15.0,0) {$\cdots$};

\draw[arr] (s0) -- node[above]{$\lambda$} (s1);
\draw[arr] (s1) -- node[above]{$\lambda$} (s2);
\draw[arr] (s2) -- node[above]{$\lambda$} (6.2,0);
\draw[arr] (6.2,0) -- node[above]{$\lambda$} (sk1);
\draw[arr] (sk1) -- node[above]{$\lambda$} (sk);
\draw[arr] (sk) -- node[above]{$\lambda$} (sk2);
\draw[arr] (sk2) -- node[above]{$\lambda$} (15.0,0);

\draw[dep] (s1) to[bend left=50] node[below]{$\mu(1)$} (s0);
\draw[dep] (s2) to[bend left=50] node[below]{$2\,\mu(2)$} (s1);
\draw[dep] (sk1) to[bend left=50] node[below]{$(k{-}1)\mu(k{-}1)$} (6.2,-0.9);
\draw[dep] (sk) to[bend left=50] node[below]{$k\,\mu(k)$} (sk1);
\draw[dep] (sk2) to[bend left=50] node[below]{$k\,\mu(k)$} (sk);
\draw[dep] (15.0,-0.9) to[bend left=50] node[below]{$k\,\mu(k)$} (sk2);

\node[align=left, anchor=north west, font=\scriptsize] at (0,-2.0) {
  \textcolor{green!35!black}{\rule{2mm}{2mm}} occupancy-зависимая зона ($j<k$): интенсивность смерти меняется с $j$\\[1pt]
  \textcolor{blue!35!black}{\rule{2mm}{2mm}} зона с прижатой интенсивностью ($j\ge k$): интенсивность смерти постоянна, $k\,\mu(k)$ -- отсюда точная смесь эрлангов для $W$
};
\end{tikzpicture}}

\captionof{figure}{Процесс гибели-размножения по числу заявок в системе $j$: при $j\ge k$ интенсивность смерти не зависит от $j$, поскольку занятость жёстко зафиксирована на потолке $k$}
\label{fig:birth-death}
\end{center}

\noindent\textbf{Регрессионные тождества.} При $\mu(n)\equiv\mathrm{const}$ для всех $n$ результат должен точно совпасть с классической системой $M/M/k/N$; при $k=1$ -- с классической $M/M/1/N$ при любой зависимости (единственное occupancy-значение $n=1$). Оба тождества подтверждены точно (таблица~\ref{tab:regression}).

Третий, качественный (не численно проверяемый ввиду вырожденности предела) вырожденный случай -- $k\to\infty$ с $\mu(n)=C/n$: занятость в этом пределе никогда не достигает потолка, $W\equiv 0$, и остаётся только разделение единой производительности $C$ поровну между всеми находящимися в системе заявками -- это ровно дисциплина egalitarian processor sharing (EPS) \cite{yashkov1987}; её перенос на $n$-канальный случай соответствующим пересчётом производительности изложен в \cite{ryzhikov2018}. Принципиальное отличие настоящей постановки от EPS -- именно в невырожденном случае конечного $k$: EPS по построению не знает понятия отказа в немедленном допуске (мощность делится между сколь угодно большим числом заявок), поэтому у неё нет и не может быть содержательного понятия времени ожидания до начала обслуживания -- вся соответствующая литература говорит о времени пребывания, не о времени ожидания в очереди. Конечный потолок занятости $k$ и FCFS-очередь вне его -- именно то структурное отличие, которое впервые делает время ожидания $W$ ненулевым и нетривиальным для этого класса систем.

\begin{center}
\captionof{table}{Точная регрессия к классической $M/M/k/N$ при occupancy-независимой интенсивности}
\label{tab:regression}

\begin{tabular}{|c|c|c|c|c|}
\hline
$k$ & $\lambda$ & $\mu$ & $\mathrm{E}[W]$ (данная модель) & $\mathrm{E}[W]$ ($M/M/k/N$) \\
\hline
4  & 3,0  & 1,00 & 0,509434 & 0,509434 \\
8  & 6,0  & 1,00 & 0,178491 & 0,178491 \\
16 & 10,0 & 0,80 & 0,094323 & 0,094323 \\
\hline
\end{tabular}
\end{center}

\articlesection{Численные эксперименты}

\noindent\textbf{Проверка независимым имитационным моделированием.} Для occupancy-зависимой интенсивности $\mu(n)=1/(\tau_0+\tau_1 n)$ (иллюстративные параметры $\tau_0=0{,}05$, $\tau_1=0{,}01$ -- структура <<базовый расход на итерацию + расход на занятое место>>, в духе параметризации \cite{ramani2026}, но не повторяющая её конкретные подогнанные значения) и потолка $k=4$ точный расчёт сопоставлен с независимым (написанным с нуля, не переиспользующим библиотечный код) дискретно-событийным имитационным моделированием при трёх уровнях загрузки (таблица~\ref{tab:des}, $3\times10^6$ заявок на прогон).

\begin{center}
\captionof{table}{Расчёт и имитация ($k=4$, occupancy-зависимая интенсивность, $3\times10^6$ заявок на прогон)}
\label{tab:des}

\begin{tabular}{|c|c|c|c|c|c|c|}
\hline
\multirow{2}{*}{$\rho$} & \multicolumn{3}{c|}{$\mathrm{E}[W]$} & \multicolumn{3}{c|}{$P(W>0{,}05)$} \\
\cline{2-7}
 & теор. & ИМ & отн. ошибка & теор. & ИМ & отн. ошибка \\
\hline
0,3 & 0,00077 & 0,00078 & 0,58\,\% & 0,00508 & 0,00514 & 1,06\,\% \\
0,5 & 0,00601 & 0,00605 & 0,64\,\% & 0,04398 & 0,04421 & 0,52\,\% \\
0,7 & 0,02805 & 0,02824 & 0,68\,\% & 0,19202 & 0,19279 & 0,40\,\% \\
\hline
\end{tabular}
\end{center}

\noindent\textbf{Эффект игнорирования occupancy-зависимости.} Естественное инженерное упрощение -- измерить интенсивность при полной занятости $\mu(k)$ и использовать её как константу, то есть свести систему к классической $M/M/k/N$ с единственной интенсивностью $\mu(k)$ (приближение в духе \cite{ramani2026}, но не воспроизводящее их конкретный метод). Таблица~\ref{tab:approx-vs-k} сравнивает точный результат с этим приближением при фиксированной относительной загрузке $\lambda=0{,}6\cdot k\cdot\mu(k)$ и растущем потолке $k$ (та же occupancy-зависимая интенсивность, что выше).

\begin{center}
\captionof{table}{Точное $\mathrm{E}[W]$ против приближения по интенсивности полной занятости, при фиксированной относительной загрузке $0{,}6$}
\label{tab:approx-vs-k}

\small
\begin{tabular}{|c|c|c|c|c|}
\hline
\multirow{2}{*}{$k$} & \multirow{2}{*}{$\mathrm{E}[W]$ точное} & \multirow{2}{*}{$\mathrm{E}[W]$ приближение} & \multicolumn{2}{c|}{Во сколько раз больше} \\
\cline{4-5}
 & & & $\mathrm{E}[W]$ & $P(W>1)$ \\
\hline
4  & 0,01328 & 0,01615  & 1,22  & 1,22 \\
8  & 0,00282 & 0,00567  & 2,01  & 2,01 \\
16 & 0,00016 & 0,00137  & 8,84  & 11,3 \\
32 & \num{4e-7} & \num{1.5e-4} & 491 & \num{9.5e3} \\
\hline
\end{tabular}
\end{center}

Приближение систематически завышает время ожидания и вероятность нарушения дедлайна -- оно приписывает всем состояниям медленную интенсивность полной занятости, тогда как при малой занятости (вдали от потолка) реальная интенсивность заметно выше. Расхождение растёт с $k$: при $k=4$ приближение завышает умеренно (в 1,2 раза), при $k=32$ -- на три-четыре порядка. Это важно практически: чем больше памяти ускорителя отведено под KV-cache (то есть чем больше допустимый потолок занятости $k$), тем дороже обходится инженерное упрощение, игнорирующее occupancy-зависимость интенсивности.

\begin{center}
\captionof{table}{Точный хвост против приближения по интенсивности полной занятости, по сетке дедлайнов ($k=16$, та же относительная загрузка $0{,}6$)}
\label{tab:deadline-grid}

\begin{tabular}{|c|c|c|c|}
\hline
$D$ & $P(W>D)$ точная & $P(W>D)$ приближение & Во сколько раз больше \\
\hline
0,02 & \num{2.57e-3} & \num{2.28e-2} & 8,84 \\
0,05 & \num{1.03e-3} & \num{9.12e-3} & 8,84 \\
0,10 & \num{2.25e-4} & \num{1.99e-3} & 8,84 \\
0,20 & \num{1.07e-5} & \num{9.44e-5} & 8,84 \\
0,40 & \num{2.40e-8} & \num{2.13e-7} & 8,84 \\
\hline
\end{tabular}
\end{center}

При фиксированном $k$ (в отличие от роста по $k$ в таблице~\ref{tab:approx-vs-k}) кратность расхождения в этом примере практически не меняется по всей проверенной сетке дедлайнов -- приближение искажает общую вероятностную массу хвоста систематически, а не его форму на этом участке.

\articlesection{Ограничения}

Фиксированный, не зависящий от истории заявки объём ресурса на одно место ($k$ -- экзогенная константа): реальный KV-cache растёт по мере генерации токенов конкретным запросом, то есть фактический потолок одновременно допустимых запросов зависит от состава уже резидентных, а не только от их числа. Обобщение на растущий во времени ресурсный след -- отдельное направление (ближайшая существующая структура в библиотеке most-queue -- модели multiserver-job).

Остаточное время обслуживания предполагается экспоненциальным (марковским) при каждом уровне occupancy; реальная оставшаяся длина генерации, вообще говоря, не экспоненциальна. Перенос уже имеющегося в библиотеке аппарата фазового (Эрланга/$H_2$) обслуживания на occupancy-зависимую интенсивность -- естественное расширение, не рассмотренное здесь.

Не рассмотрены: двухфазность обслуживания (отдельные по характеру фазы prefill и decode современных инференс-движков); нетерпение заявок; само время обслуживания после допуска в активный пул (величина $W$ в данной работе -- только задержка до допуска, не полное время пребывания в системе).

Численное сравнение с <<приближением по интенсивности полной занятости>> -- собственная, явно обозначенная как иллюстративная конструкция автора, в духе общей структуры приближённого метода \cite{ramani2026}, но не повторяющая его конкретные подогнанные параметры и не являющаяся прямым воспроизведением чьих-либо опубликованных численных результатов.

\articlesection{Заключение}

Получено и верифицировано точное распределение времени ожидания (начальные моменты и хвост) для occupancy-зависимой системы массового обслуживания с потолком занятости произвольной величины $k$ и произвольным профилем зависимости интенсивности от занятости -- марковской модели ядра continuous batching современных LLM-инференс-движков. Результат обобщает точный одно-/двухсерверный результат Morrison (1989) на произвольное $k$ и даёт, в отличие от единственной найденной близкой по постановке современной работы, точный (не приближённый методом среднего значения) результат, включая точный хвост, а не только среднее. Численно показано, что инженерное упрощение, игнорирующее occupancy-зависимость интенсивности, систематически завышает оценку риска нарушения дедлайна, причём расхождение растёт с ростом допустимого потолка занятости -- на один-два порядка при умеренном $k$ и на несколько порядков при большом. Открытая программная реализация доступна в библиотеке most-queue \cite{mostqueue2026}.

\articlesection{Список литературы}
\printbibliography[heading=none]

\end{document}
